\documentclass[journal]{IEEEtran}
\usepackage[a5paper, margin=10mm, onecolumn]{geometry}
\usepackage{lmodern}
\usepackage{tfrupee}
\setlength{\headheight}{1cm}
\setlength{\headsep}{0mm}

\usepackage{gvv-book}
\usepackage{gvv}
\usepackage{cite}
\usepackage{amsmath,amssymb,amsfonts,amsthm}
\usepackage{algorithmic}
\usepackage{graphicx}
\usepackage{textcomp}
\usepackage{xcolor}
\usepackage{txfonts}
\usepackage{listings}
\usepackage{enumitem}
\usepackage{mathtools}
\usepackage{gensymb}
\usepackage{comment}
\usepackage[breaklinks=true]{hyperref}
\usepackage{tkz-euclide}
\usepackage{listings}
\def\inputGnumericTable{}
\usepackage[latin1]{inputenc}
\usepackage{color}
\usepackage{array}
\usepackage{longtable}
\usepackage{calc}
\usepackage{multirow}
\usepackage{hhline}
\usepackage{ifthen}
\usepackage{lscape}
\usepackage{xparse}

\bibliographystyle{IEEEtran}

\title{Software Project - Report}
\author{EE25BTECH11043 - Nishid Khandagre}

\begin{document}
\maketitle

\renewcommand{\thefigure}{\theenumi}
\renewcommand{\thetable}{\theenumi}

\textbf{Summary of Strang's Video}


Every matrix $A$ can be expressed as:
\[
A = U \Sigma V^T
\]
where $U$ and $V$ are orthogonal matrices whose columns are orthonormal vectors,\\
$\Sigma$ is a diagonal matrix with entries $\sigma_i$ called the \textbf{singular values}.

\[
A^T A = V (\Sigma^T \Sigma) V^T
\]
\[
A A^T = U (\Sigma \Sigma^T) U^T
\]


\textbf{Properties}
\begin{itemize}
    \item $A v_i = \sigma_i u_i$
    \item The column vectors $v_i$ are eigenvectors of $A^T A$:
    \item The column vectors $u_i$ are eigenvectors of $A A^T$ :
\end{itemize}

\textbf{Algorithm Used: Jacobian}


The Jacobian method is an iterative algorithm for computing the Singular Value Decomposition (SVD) of a matrix $\vec{A}$.\\
It applies rotations to eliminate off-diagonal elements of \(A^T A \), and converting it into a diagonal matrix. The diagonal elements of the obtained diagonal matrix represent the square of the singular values of \( A \).


\[
A=U \Sigma V^T
\]
\[
A^T A = V \Sigma^T\Sigma V^T
\]
\[
U = A V \Sigma^{-1}
\]
Here, \( V \) and \( U \) are orthogonal matrices containing right and left singular vectors, and \( \Sigma \) is the diagonal matrix of singular values.

\textbf{Explanation}

\begin{enumerate}
\item \textbf{Form the symmetric matrix}

\[
M = A^T A
\] 
    \item \textbf{Find the largest off-diagonal element}  \\
    Identify indices \( (p, q) \) for the largest off-diagonal element where \( p < q \).
    
    Using this element we will find the angle.

    \item \textbf{Calculate the rotation angle}\\
    The rotation angle \( \theta \) is chosen such that the new matrix has
    a zero at position \( (p, q) \)
    \[
    \tan(2\theta) = \frac{2M_{pq}}{M_{qq} - M_{pp}}
    \]
    Then
    \[
    c = \cos(\theta), \qquad s = \sin(\theta)
    \]

    \item \textbf{Form the orthogonal rotation matrix \( S \)} \\ 
    Construct an identity matrix and modify the four elements
    corresponding to indices \( (p, q) \):
    \[
    S_{pp} = S_{qq} = c , \quad S_{pq} = s, \quad S_{qp} = -s
    \]
    All other diagonal entries are 1 and the rest are 0.

    \item \textbf{Find the transformation matrix} 
    \[
    R = S^T M S
    \]
    By This operation the off-diagonal term \( M_{pq} \) placed at (p,q) becomes zero in $R$.\\

    \item \textbf{Iterate}  \\
    Repeat steps 1 to 4 for different pairs \( (p, q) \) until all off-diagonal elements become small or zero 

    \item \textbf{Obtain eigenvalues}  
    When the algorithm is done , \( R \) becomes approximately diagonal.
    where each diagonal entry \( \lambda_i = R_{ii} \) is an eigenvalue of the original matrix.

    \item \textbf{Compute eigenvectors} \\ 
    The final orthogonal transformation matrix is obtained by
    successively multiplying all rotation matrices used:
    \[
    V = S_1 S_2 S_3 \cdots S_k
    \]
    Each time we multiply the rotation orthogonal matrix with V
   
    The columns of \( V \) are the normalized eigenvectors corresponding
    to the eigenvalues \( \lambda_i \).
\end{enumerate}


\textbf{Pseudocode of the Algorithm}
\begin{enumerate}
    \item \textbf{Read Input Image:}  
    Read the grayscale image and store it as matrix \( A \) and also normalize each entry of A .  
    
    Extract number of rows \( r \), columns \( c \), and maximum pixel value .

    \item \textbf{Compute Transpose:}  
    Compute the transpose of the image matrix:

    \item \textbf{Form Symmetric Matrix:}  
    Construct the symmetric matrix \( A^T A \):
    \[
    B = A^T A
    \]

    \item \textbf{Initialize Arrays that will stor the eigen vectors:}  
    \[
    \vec{V} = I
    \]
   
    \item \textbf{Perform Jacobi Eigenvalue Decomposition:}\\ 
    Repeat until convergence or until the iteration limit is reached:
    \begin{enumerate}
        \item Find indices of maximum off diagonal element \( (p, q) \)
        \item If max off-diagonal element is very small , stop iteration.
        \item Compute rotation angle:
        \[
        \theta = \frac{1}{2} \tan^{-1}\left( \frac{2 B_{pq}}{B_{qq} - B_{pp}} \right)
        \]
        \item Compute cosine and sine of the angle:
        \[
        c = \cos(\theta), \quad s = \sin(\theta)
        \]
        \item Form the rotation orthogonal matrix S and then find the transformation matrix 
        \[
        B = S^T B S
        \]
        \item Update eigenvector matrix:
        \[
        \vec{V} = \vec{V} S
        \]
    \end{enumerate}
    After convergence, the diagonal elements of \( B \) contain eigenvalues:

    \item \textbf{Sort Eigenvalues and Vectors:}  
    Sort eigenvalues \( \lambda_i \) in descending order and reorder columns of \( \vec{V} \) accordingly.

    \item \textbf{Compute Singular Values:}  
    The singular values are obtained as:
    \[
    \sigma_i = \sqrt{\lambda_i}
    \]
    forming the diagonal matrix:
    \[
    \Sigma = \text{diag}(\sigma_1, \sigma_2, \dots, \sigma_c)
    \]

    \item \textbf{Compute Left Singular Vectors:}  
    Each left singular vector is computed as:
    \[
    \vec{u}_j = \frac{A \vec{v}_j}{\sigma_j}
    \]

    \item \textbf{Construct top k Matrix:}  
    Using the top \( k \) singular values and vectors:
    \[
    A_k = \sum_{t=1}^{k} \sigma_t \, \vec{u}_t \vec{v}_t^T
    \]

    \item \textbf{Calculate Frobenius Error:}  

    \[
    F = \sqrt{ \sum_{i=1}^{r} \sum_{j=1}^{c} (A_{ij} - A_{ij}^{k})^2 }
    \]

    \item \textbf{Save Compressed Image:}  
    Save the reconstructed image matrix \\

\end{enumerate}

\textbf{Comparison of Different Algorithms}
\begin{enumerate}
    \item \textbf{Householder Bidiagonalization and QR Algorithm} \\   
    The algorithm first reduces the input matrix to a bidiagonal form using Householder reflections, 
    and then applies iterative QR transformations to extract singular values and vectors.  
    It is highly efficient but it is complex to implement manually.\\

    \item \textbf{Power Iteration}  \\
    These methods compute one eigenvalue and eigenvector at a time by iteratively applying 
    the matrix to a random vector until convergence.  
    power iteration is computationally slow and can miss smaller eigenvalues.\\

    \item \textbf{Jacobian Method}  \\
    The Jacobi algorithm applies a series of orthogonal rotations 
    to eliminate off-diagonal elements of a symmetric matrix.  
    It is simple to understand and implement, guarantees numerical stability, 
    and always produces an orthogonal eigenvector matrix.\\
\end{enumerate}

\textbf{Why Jacobi Method was Chosen}\\
\begin{itemize}
    \item it is easy to implement, and conceptually easy. 
    \item It provides a clear understanding of how eigenvalues and eigenvectors are obtained through orthogonal transformations.  
    \item Each step involves simple arithmetic and trigonometric operations, making it simple to understand.
\end{itemize}

\newpage

\textbf{Reconstructed images}\\\\

\textbf{For einstein.pgm} - Original image
\begin{figure}[H]
        \centering
        \includegraphics[width=0.5\columnwidth]{einstein.jpg}
        \caption{}
        \label{fig:q39}
    \end{figure}
 \begin{enumerate}
        \item For k=20
        \begin{figure}[H]
        \centering
        \includegraphics[width=0.5\columnwidth]{einstein20.png}
        \caption{}
        \label{fig:q39}
    \end{figure}
    \item For k=50
    \begin{figure}[H]
        \centering
        \includegraphics[width=0.5\columnwidth]{einstein50.png}
        \caption{}
        \label{fig:q39}
    \end{figure}
    \item For k=100
    \begin{figure}[H]
        \centering
        \includegraphics[width=0.5\columnwidth]{einstein100.png}
        \caption{}
        \label{fig:q39}
    \end{figure}
    \item For k=180
    \begin{figure}[H]
        \centering
        \includegraphics[width=0.5\columnwidth]{einstein180.png}
        \caption{}
        \label{fig:q39}
    \end{figure}    
    \end{enumerate}

\textbf{For globe.pgm} - Original image
\begin{figure}[H]
        \centering
        \includegraphics[width=0.5\columnwidth]{globe.jpg}
        \caption{}
        \label{fig:q39}
    \end{figure}
 \begin{enumerate}
        \item For k=200
        \begin{figure}[H]
        \centering
        
\includegraphics[width=0.5\columnwidth]{globe200.png}
        \caption{}
    \end{figure}
    \item For k=400
    \begin{figure}[H]
        \centering
        \includegraphics[width=0.5\columnwidth]{globe400.png}
        \caption{}
        \label{fig:q39}
    \end{figure}
    \item For k=600
    \begin{figure}[H]
        \centering
        \includegraphics[width=0.5\columnwidth]{globe600.png}
        \caption{}
        \label{fig:q39}
    \end{figure}
    \item For k=800
    \begin{figure}[H]
        \centering
        \includegraphics[width=0.5\columnwidth]{globe800.png}
        \caption{}
        \label{fig:q39}
    \end{figure}    
    \end{enumerate}

\textbf{For greyscale.pgm} - Original image
\begin{figure}[H]
        \centering
        \includegraphics[width=0.5\columnwidth]{greyscale.png}
        \caption{}
        \label{fig:q39}
    \end{figure}
 \begin{enumerate}
        \item For k=400
        \begin{figure}[H]
        \centering
        \includegraphics[width=0.5\columnwidth]{greyscale400.png}
        \caption{}
    \end{figure}
    \item For k=600
    \begin{figure}[H]
        \centering
        \includegraphics[width=0.5\columnwidth]{greyscale600.png}
        \caption{}
        \label{fig:q39}
    \end{figure}
    \item For k=800
    \begin{figure}[H]
        \centering
        \includegraphics[width=0.5\columnwidth]{greyscale800.png}
        \caption{}
        \label{fig:q39}
    \end{figure}
    \item For k=1000
    \begin{figure}[H]
        \centering
        \includegraphics[width=0.5\columnwidth]{greyscale1000.png}
        \caption{}
        \label{fig:q39}
    \end{figure}    
    \end{enumerate}



\textbf{Error Analysis}
The error was calculated using the Frobenius norm:
\[
E_k = \| A - A_k \|_F = \sqrt{ \sum_{i,j} (A_{ij} - A_{ij}^{k})^2 }
\]

\begin{enumerate}

\item For einstein.pgm
\begin{table}[H]
    \centering
    \begin{tabular}{|c|c|}
        \hline
        \textbf{k value} & \textbf{Forbenius Error} \\  
        \hline
        5 & 20201.95 \\                              
        \hline
        20 & 18543.28 \\                         
        \hline
        50 & 15997.30 \\       
        \hline
        100 & 11903.16 \\       
        \hline
    \end{tabular}
       \label{tab:1}
\end{table}

\item For globe.pgm
\begin{table}[H]
    \centering
    \begin{tabular}{|c|c|}
        \hline
        \textbf{k value} & \textbf{Forbenius Error} \\  
        \hline
        5 & 154428.09 \\                              
        \hline
        20 & 152464.03 \\                         
        \hline
        50 & 148137.73 \\       
        \hline
        100 & 140759.91 \\       
        \hline
    \end{tabular}
    \label{tab:2}
\end{table}

\item For greyscale.pgm
\begin{table}[H]
    \centering
    \begin{tabular}{|c|c|}
        \hline
        \textbf{k value} & \textbf{Forbenius Error} \\  
     \hline
        5 & 190586.46 \\                              
        \hline
        20 & 188237.54 \\                         
        \hline
        100 & 173920.33 \\       
        \hline
        200 & 157032.58 \\       
        \hline
    \end{tabular}
    \label{tab:3}
\end{table}

\end{enumerate}


\textbf{Discussion and Trade-off}\\
The $k$ controls how many singular values are retained during image reconstruction.  
A smaller \(k\) gives higher compression but results in loss of fine details and a blurrier image.  
A larger \(k\) preserves more detail and improves image quality, but requires more storage.  
Thus, \(k\) balances between quality and compression efficiency.\\

\textbf{Conclusion}\\
The project successfully implemented image compression using the truncated SVD with the Jacobian method.  

\end{document}